\begin{tcolorbox}[enhanced,colback=paper,colframe=linegray,boxrule=.7pt,arc=2mm,
left=3mm,right=3mm,top=3mm,bottom=2mm,title={图解：Query 与三把 Key 匹配，Softmax 后读取 Value},
fonttitle=\sffamily\bfseries,coltitle=ink]
\centering
\begin{tikzpicture}[>=Latex,
 block/.style={draw=linegray,fill=white,rounded corners=2pt,minimum width=27mm,
 minimum height=9mm,align=center,font=\small},
 step/.style={draw=blue,fill=blueback,rounded corners=2pt,
 minimum width=28mm,minimum height=10mm,align=center,font=\small}]
\node[step] (q) at (0,0) {Query\\$\vect q=(1,0)$};
\node[block] (k1) at (4,1.35) {$k_1=(2,0)$};
\node[block] (k2) at (4,0) {$k_2=(0,1)$};
\node[block] (k3) at (4,-1.35) {$k_3=(1,1)$};
\node[step] (s1) at (8,1.35) {$s_1=2$};
\node[step] (s2) at (8,0) {$s_2=0$};
\node[step] (s3) at (8,-1.35) {$s_3=1$};
\foreach \k/\s in {k1/s1,k2/s2,k3/s3}
 \draw[->,thick,blue] (\k)--node[above,font=\scriptsize]{点积}(\s);
\foreach \k in {k1,k2,k3}
 \draw[->,teal] (q.east) -- (\k.west);
\node[draw=teal,fill=greenback,rounded corners=3pt,align=center,
 minimum width=93mm,minimum height=11mm,font=\small] (w) at (4,-3.15)
 {Softmax：$(2,0,1)\ \longrightarrow\ (0.665,\ 0.090,\ 0.245)$};
\draw[->,orange,thick] (s3.south) to[out=-90,in=0] (w.east);
\node[draw=teal,fill=greenback,rounded corners=3pt,align=center,
 minimum width=93mm,minimum height=11mm,font=\small] (out) at (4,-4.65)
 {$0.665v_1+0.090v_2+0.245v_3=(0.910,\ 0.335)$};
\draw[->,thick,teal] (w)--node[right,font=\scriptsize]{加权读取} (out);
\end{tikzpicture}
\par\smallskip
{\small\raggedright\textbf{读图：}蓝框是匹配分数，绿色框是归一化权重和最终读出的向量。
本图使用本节的未缩放点积分数；标准缩放点积 Attention 还会除以 $\sqrt{d_k}$。\par}
\end{tcolorbox}